\begin{table}[t]
\tbl{\label{tab:orthogonal:activities}}
{\begin{tcolorbox}[tabularx={p{3cm}|p{2cm}|p{6.3cm}|},enhanced,fonttitle=\bfseries\large,fontupper=\normalsize\sffamily,
colback=yellow!10!white,colframe=red!50!black,colbacktitle=blue!30!white,
coltitle=black,title=Orthogonal Space Activities]
\hline
Command & Arguments & Purpose \\
\hline
\hline
\verb'-cheat_sheet_orthogonal'  &  options  &  Create a cheat sheet. Considers the given draw options object. \\
\hline
\verb'-print_points'  &  $v$  &  Print the points whose ranks are given in the vector $v$. \\
\hline
\verb'-print_lines'  &  $v$  &  Print the lines whose ranks are given in the vector $v$. \\
\hline
\verb'-unrank_line_through_two_points'  &  $p1$ $p2$  & Determine the rank of the line through the points whose ranks are $p1$ and $p2$. \\
\hline
\verb'-lines_on_point'  &  $p$  &   Create the ranks of all lines through the point whose rank is $p$. \\
\hline
\verb'-perp'   &  $L$  & Determine the common perp of a set of points. The point ranks are given in the list $L$.  \\
\hline
\verb'-export_point_line_incidence_matrix'   &    & Create a csv file with the point line incidence matrix of the space.  \\
\hline
\verb'-intersect_with_subspace'   &  $M$  & Find the points in the intersection of the quadric with the subspace whose generating matrix has label $M$. \\
\hline
\verb'-intersect_with_subspace_given_by_base'   &  $base$  & Find the points in the intersection of the quadric with the subspace whose base elements are given by means of their Orbiter ranks. \\
\hline
\verb'-table_of_blt_sets'   &   &  Export a table of BLT-sets. We must be in an $Q^4(q)$ space.  \\
\hline
\verb'-create_orthogonal_reflection'   &  points  &  Create the orthogonal reflections associated with the given set of points, encoded as a vector. The points must be off the quadric. The points are given as point ranks in projective space. \\
\hline
\verb'-create_perp_of_points'   &  points  &  For each point p in the given set of points, compute the set of points q on the quadric which are perpendicular to p. The input set of points is coded in the underlying projective geometry. The output points are coded as points on the quadric.  \\
\hline
\verb'-create_Siegel_transformation'   &  u v  &  Create the Siegel transformation defined by the vectors $u$ and $v$. The Siegel transformation is defined as $\rho_{u,v}(x) := x + \beta(x,v) u - \beta(x,u) v - Q(v) \beta(x,u) u.$  Here, $Q$ is the quadratic form and $\beta$ is the associated bilinear form. The vector $u$ must be singular and $v$ must be in $u^\perp.$ \\
\hline
\verb'-make_all_Siegel_transformations'   &   &  Make all Siegel transformations by looping over all possible pairs for $u$ and $v$.  \\
\hline
\verb'-create_orthogonal_reflection_6_and_4'   &  points A4  &  Create the orthogonal reflections associated with the given set of point. Once done, turn the mappings into mappings of $\PG(3,q)$, using the action A4. \\
\hline
\end{tcolorbox}}
\end{table}
